\documentclass[a4paper]{article}
% Español (México) con XeLaTeX: fontspec para fuentes Unicode, polyglossia
% para la división silábica y los nombres del documento en la variante mexicana.
\usepackage{fontspec}
\usepackage{polyglossia}
\setdefaultlanguage[variant=mexican]{spanish}
% Los nombres chinos van como «pinyin (original)»: xeCJK da a los caracteres
% Han su propia fuente; sin ella XeLaTeX los omite con un aviso.
% FandolSong no cubre algunos caracteres de nombres (U+73FA); WenQuanYi Zen Hei sí.
\usepackage[AutoFallBack=true]{xeCJK}
\setCJKmainfont{FandolSong-Regular.otf}
\setCJKfallbackfamilyfont{\CJKrmdefault}{WenQuanYi Zen Hei}
% polyglossia no traduce los nombres de listings.
\AtBeginDocument{\providecommand{\lstlistingname}{}\renewcommand{\lstlistingname}{Listado}%
  \providecommand{\lstlistlistingname}{}\renewcommand{\lstlistlistingname}{Índice de listados}}
\usepackage{amsmath,amssymb}
\usepackage{graphicx}
\usepackage[margin=2.35cm]{geometry}
\usepackage[most]{tcolorbox}
\usepackage{etoolbox}
\usepackage{listings}
\usepackage{xcolor}
\usepackage{hyperref}
\usepackage{booktabs,longtable,tabularx,array}
\usepackage{subcaption}
\usepackage{float}
\usepackage{enumitem}
\usepackage{microtype}
\usepackage{fancyhdr}
\usepackage{needspace}

\definecolor{kbBack}{HTML}{F2F6FA}
\definecolor{kbFrame}{HTML}{9DB4CC}
\definecolor{kbTitle}{HTML}{52708F}
\definecolor{ibBack}{HTML}{FBF5E7}
\definecolor{ibFrame}{HTML}{C9A961}
\definecolor{ibTitle}{HTML}{7D5F1E}
\definecolor{wbBack}{HTML}{FCF2F0}
\definecolor{wbFrame}{HTML}{D6A096}
\definecolor{wbTitle}{HTML}{9E4F43}
\definecolor{dbBack}{HTML}{F4F7F3}
\definecolor{dbFrame}{HTML}{AEC2AC}
\definecolor{dbTitle}{HTML}{5F7F64}

\tcbset{softbox/.style={
  enhanced,breakable,boxrule=0.8pt,arc=2pt,left=3mm,right=3mm,
  top=2mm,bottom=2mm,coltitle=white,fonttitle=\bfseries,
  toptitle=0.6mm,bottomtitle=0.6mm
}}
\newtcolorbox{knowledgebox}[1]{softbox,colback=kbBack,colframe=kbFrame,colbacktitle=kbTitle,title={#1}}
\newtcolorbox{importantbox}[1]{softbox,colback=ibBack,colframe=ibFrame,colbacktitle=ibTitle,title={#1}}
\newtcolorbox{warningbox}[1]{softbox,colback=wbBack,colframe=wbFrame,colbacktitle=wbTitle,title={#1}}
\newtcolorbox{dialoguebox}[1]{softbox,colback=dbBack,colframe=dbFrame,colbacktitle=dbTitle,title={#1}}

\lstset{
  language=Python,basicstyle=\ttfamily\small,
  keywordstyle=\color{kbTitle}\bfseries,stringstyle=\color{wbTitle},
  commentstyle=\color{dbTitle},breaklines=true,frame=single,numbers=left,
  numberstyle=\tiny\color{gray},captionpos=b,columns=fullflexible,
  keepspaces=true,showstringspaces=false,extendedchars=true
}
\BeforeBeginEnvironment{lstlisting}{\Needspace{12\baselineskip}}

\hypersetup{colorlinks=true,linkcolor=kbTitle,urlcolor=kbTitle,citecolor=wbTitle}
\setlist{itemsep=0.25em,topsep=0.4em}
\setlength{\parindent}{2em}
\setlength{\parskip}{0.25em}
\newcolumntype{L}[1]{>{\raggedright\arraybackslash}p{#1}}

\providecommand{\notetitle}{Notas del curso: [tema del curso]}
\providecommand{\noteauthors}{Elaboradas a partir de los materiales públicos de Stanford CS25}
\providecommand{\notedate}{Versión reescrita en 2026}
\providecommand{\videochannel}{Stanford Online}
\providecommand{\videopublishdate}{2022}
\providecommand{\videoduration}{[duración del video]}
\providecommand{\videourl}{}
\providecommand{\videocoverpath}{cover.jpg}
\providecommand{\coursesite}{https://web.stanford.edu/class/cs25/}
\providecommand{\repourl}{https://github.com/hqhq1025/ai-course-notes}
\providecommand{\skillversion}{wdkns/wdkns-skills@39f1a04}

\newcommand{\figsource}[1]{\par\vspace{-0.3em}{\footnotesize\color{gray}\emph{Fuente: #1}}\par}
\newcommand{\teachervoice}[1]{\begin{knowledgebox}{Nota de clase}#1\end{knowledgebox}}
\newcommand{\term}[1]{\textbf{#1}}

\pagestyle{fancy}
\fancyhf{}
\fancyfoot[C]{\thepage}
\fancyfoot[R]{\footnotesize\color{gray}\href{\repourl}{AI Course Notes}}
\renewcommand{\headrulewidth}{0pt}
\renewcommand{\footrulewidth}{0pt}

\newcommand{\makecscover}{
\begin{titlepage}
\centering
\vspace*{0.7cm}
{\Large Stanford CS25 · Transformers United\par}
\vspace{0.8cm}
{\huge\bfseries \notetitle\par}
\vspace{0.6cm}
{\large \noteauthors\par}
\vspace{0.25cm}
{\large \notedate\par}
\vspace{0.9cm}
\ifdefempty{\videocoverpath}{}{
  \includegraphics[width=0.86\textwidth,height=0.43\textheight,keepaspectratio]{\videocoverpath}\par
}
\vfill
\begin{tcolorbox}[width=0.92\textwidth,colback=black!2!white,colframe=black!45,boxrule=0.7pt,arc=2pt]
\textbf{Autor o canal del video}: \videochannel\par
\textbf{Fecha de publicación}: \videopublishdate\par
\textbf{Duración del video}: \videoduration\par
\textbf{Sitio del curso}: \href{\coursesite}{\nolinkurl{\coursesite}}\par
\ifdefempty{\videourl}{}{\textbf{Enlace del video}: \href{\videourl}{\nolinkurl{\videourl}}\par}
\textbf{Normas de elaboración}: \skillversion\ + las normas source-first / teacher-voice / visual-QA de este repositorio
\end{tcolorbox}
\end{titlepage}
\tableofcontents
\newpage
}

\usepackage{multicol}

\renewcommand{\notetitle}{Lecture 27: de la Scaling Intuition a la evolución de la arquitectura Transformer}
\renewcommand{\noteauthors}{Elaborado a partir de la clase de Jason Wei y Hyung Won Chung, dos deck oficiales y subtítulos hechos a mano}
\renewcommand{\notedate}{CS25 V4 · versión reescrita source-first de 2026}
\renewcommand{\videochannel}{Stanford Online}
\renewcommand{\videopublishdate}{Clase: 2024-04-11; subido: 2024-05-06}
\renewcommand{\videoduration}{1:17:07}
\renewcommand{\videourl}{https://www.youtube.com/watch?v=3gb-ZkVRemQ}
\renewcommand{\videocoverpath}{cover.jpg}

\newcommand{\lecturefigure}[3]{%
\begin{figure}[H]
  \centering
  \includegraphics[width=0.97\textwidth,height=0.49\textheight,keepaspectratio]{slides-images/#1}
  \caption{#2}
  \figsource{#3}
\end{figure}
}

\let\standardtableofcontents\tableofcontents
\renewcommand{\tableofcontents}{%
  \begingroup
  \small
  \setlength{\parskip}{0pt}
  \standardtableofcontents
  \endgroup
}

\begin{document}
\makecscover

\section{Dos conferencias, un método de investigación: primero mirar los datos, después estudiar el cambio}

Esta clase se compone de dos conferencias de estilo distinto y lógica complementaria. Jason Wei parte de los datos de entrenamiento y de las curvas de comportamiento para preguntarse por qué un next-token objective tan simple produce tantas capacidades; Hyung Won Chung, en cambio, parte de los cambios históricos del Transformer para preguntarse qué estructuras son atajos en las condiciones de cómputo actuales y cuáles se convertirán en cuellos de botella cuando la escala crezca. Al poner juntas las dos conferencias, no se obtienen dos grupos de opiniones dispersas, sino un método común: \textbf{primero entrar en contacto con evidencia real, después descomponer la observación en mecanismos verificables y, por último, estudiar las tendencias a lo largo del eje de la escala o del tiempo.}

\subsection{La pregunta de Jason: ¿por qué los modelos de lenguaje funcionan tan bien?}

Esta sección establece primero la postura de investigación de la primera conferencia. La portada no promete una teoría completa, sino que dice explícitamente que ofrecerá algunas intuitions. Aquí, una intuition no es una conjetura arbitraria, sino una hipótesis de trabajo condensada a partir de los datos, de las curvas de pérdida y de prompts concretos; debe poder orientar el siguiente experimento, y no solo explicar el éxito a posteriori.

\lecturefigure{jason-slide-001.jpg}{Tema de la conferencia de Jason Wei: Intuitions on Language Models}{Jason Wei official deck, p.~1}

\begin{importantbox}{Las cuatro preguntas de la primera conferencia}
\begin{enumerate}
  \item ¿Por qué la next-word prediction no es solo una tarea, sino algo parecido a una mezcla de una enorme cantidad de tareas?
  \item ¿Por qué aumentar la escala reduce de forma estable la loss global?
  \item Si la loss global es suave, ¿por qué una tarea individual presenta cambios abruptos?
  \item ¿Por qué ciertas tareas cuidadosamente seleccionadas muestran inverse o U-shaped scaling?
\end{enumerate}
\end{importantbox}

\lecturefigure{jason-slide-002.jpg}{La pregunta básica y el método de inspeccionar manualmente los datos de entrenamiento}{Jason Wei official deck, p.~2}

Al leer esta diapositiva, conviene mirar primero el método de la esquina inferior derecha, y no solo el título. ``Why do LLMs work?'' es demasiado amplia para experimentar con ella directamente; ``manually inspect data'', en cambio, convierte la pregunta en una observación ejecutable: tomar muestras aleatorias de texto, hacer la continuation a mano, marcar el conocimiento y las reglas que requiere un token y preguntarse después de qué contextos pudo haber aprendido el modelo esas capacidades. Esto vuelve a conectar la investigación abstracta sobre modelos con la distribución de entrenamiento.

\teachervoice{Entre 00:00:55 y 00:01:27, Jason llama a «mirar los datos a mano» una herramienta extremadamente útil. Esta advertencia se opone a mirar solo las puntuaciones agregadas de los benchmarks: si el investigador nunca ha leído las muestras de entrenamiento, es fácil que confunda una señal de supervisión evidente en los datos con una emergencia misteriosa.}

\lecturefigure{jason-slide-003.jpg}{Mirar los datos a mano es entrenar la propia «red neuronal biológica» del investigador}{Jason Wei official deck, p.~3}

El sentido didáctico de esta frase es que el investigador también necesita representation learning. Leer muestras una y otra vez, predecir etiquetas y revisar los errores con expertos del dominio forma en la mente representaciones del ruido, la ambigüedad, los atajos y la cola larga. Jason usa su experiencia con la clasificación de cáncer de pulmón para mostrar que aprender primero la tarea misma le dio la intuición que más tarde le permitió plantear preguntas de investigación; esto no exige que cada investigador se convierta en un experto completo del dominio, pero sí que entienda, al menos, qué significa un error.

\begin{lstlisting}[caption={Ciclo mínimo de investigación de la manual data inspection}]
for example in sample(training_corpus):
    guess = human_predict(example.context)
    compare(guess, example.next_token)
    label_required_signal(example)
    record_noise_shortcut_ambiguity(example)
summarize_recurring_patterns()
design_one_testable_experiment()
\end{lstlisting}

\subsection{La pregunta de Hyung Won: cuando no se alcanza a seguir los avances más recientes, estudiar el cambio mismo}

La subsección anterior construyó la intuición a partir de los datos; ahora se amplía la escala temporal. El título de Hyung Won subraya ``from the history of Transformer'', pero su propósito no es hacer una crónica, sino identificar en la historia las fuerzas impulsoras que se repiten. Al leer las siguientes diapositivas, ``history'' debe entenderse como un conjunto de experimentos naturales: cuando cambian el cómputo, los datos, los objetivos y la estructura, qué métodos se conservan, cuáles se descartan y por qué.

\lecturefigure{hyung-slide-001.jpg}{Tema de la conferencia de Hyung Won Chung: moldear el futuro de la IA a partir de la historia del Transformer}{Hyung Won Chung official deck, p.~1}

\lecturefigure{hyung-slide-002.jpg}{No perseguir solo los resultados más recientes: estudiar el cambio mismo}{Hyung Won Chung official deck, p.~2}

La segunda diapositiva redefine la ansiedad de investigación como un problema de método. Perseguir cada semana los modelos nuevos proporciona una gran cantidad de hechos, pero no necesariamente una dirección; estudiar el ``change itself'' exige, en cambio, alinear los sistemas antiguos con los nuevos y encontrar las variables comunes detrás de los cambios de desempeño. El objetivo no es predecir con precisión el nombre de algún producto, sino aumentar la probabilidad de acertar al juzgar la dirección de la investigación.

\teachervoice{Entre 00:32:29 y 00:33:22, Hyung Won enfatiza que, cuando la velocidad del cambio supera la capacidad humana de seguirlo, conviene más volver a mirar lo antiguo y trazar la ruta de «cómo llegamos hasta aquí». Lo importante de la clase no es memorizar la tabla de modelos de 2024, sino practicar la abstracción de una dirección a partir de la ruta.}

\lecturefigure{hyung-slide-003.jpg}{Método de tres pasos para estudiar el cambio: identificar, comprender y extrapolar la fuerza dominante}{Hyung Won Chung official deck, p.~3}

\begin{importantbox}{Definición de trabajo de dominant driving force}
La fuerza impulsora dominante no es la única fuerza del sistema, sino la que explica la mayor parte de la direccionalidad a la escala del problema actual. El investigador puede ignorar temporalmente los factores menores a cambio de poder analizar el sistema; pero debe declarar abiertamente que se trata de una aproximación y volver a examinar los factores ignorados cuando la predicción falle.
\end{importantbox}

\lecturefigure{hyung-slide-004.jpg}{Experimento de la pluma que cae: la gravedad para ilustrar cómo ignorar las fuerzas secundarias}{Hyung Won Chung official deck, p.~4}

Para una pluma que se suelta desde el reposo, si se toma la dirección hacia abajo como positiva, cerca de la superficie y despreciando la resistencia del aire, se tiene
\[
y(t)=y_0+v_0t+\frac{1}{2}gt^2.
\]
Aquí, $y(t)$ representa la posición en el instante $t$, $y_0$ es la posición inicial, $v_0$ es la velocidad inicial y $g$ es la aceleración de la gravedad. La fórmula es útil no porque la resistencia no exista, sino porque, con la precisión que requiere esta demostración, la gravedad domina lo suficiente y ya se cuenta con un modelo confiable de ella.

\teachervoice{Entre 00:35:04 y 00:36:21, la clase menciona explícitamente la resistencia del aire y las interacciones aerodinámicas complejas, pero decide ignorarlas. Lo que enseña este ejemplo no es que «la realidad tiene una sola variable», sino «primero definir la precisión requerida y después decidir qué variables pueden aproximarse».}

\lecturefigure{hyung-slide-006.jpg}{Cuantas más fuerzas dominantes y más complejas sus interacciones, más difícil es predecir la trayectoria}{Hyung Won Chung official deck, p.~6}

\begin{knowledgebox}{Lectura de la figura: no confundir el eje horizontal con el «número total de variables»}
El eje horizontal es el \textbf{número de fuerzas dominantes}, no el número de campos observados. Un sistema puede tener miles de características y, aun así, estar gobernado por unos pocos mecanismos; a la inversa, un sistema en apariencia simple también puede ser difícil de extrapolar si tiene varias retroalimentaciones fuertemente acopladas. Esta figura respalda una estrategia de modelado: primero buscar los factores con mayor direccionalidad y después decidir si hace falta incorporar sus interacciones.
\end{knowledgebox}

\subsection{El ciclo de investigación común}

Las dos subsecciones anteriores partieron, respectivamente, de un corte transversal de los datos y de un eje histórico longitudinal; ahora pueden combinarse en un ciclo. Primero se miran los datos reales para no alejarse de la distribución; después se trazan curvas o trayectorias históricas para identificar la dirección; luego se propone una explicación mínima del mecanismo y se buscan activamente contraejemplos que la refuten. La diferencia entre los dos conferencistas está solo en la escala de observación, no en la estructura epistemológica.

\begin{table}[H]
\centering
\begin{tabularx}{\textwidth}{L{0.19\textwidth} >{\raggedright\arraybackslash}X >{\raggedright\arraybackslash}X}
\toprule
Paso & Versión de Jason & Versión de Hyung Won \\
\midrule
Contacto con la evidencia & Leer a mano muestras de entrenamiento y prompts & Comparar arquitecturas antiguas con las actuales \\
Búsqueda de estructura & Descomponer latent tasks y loss mixture & Identificar la dominant driving force \\
Construcción de curvas & Trazar scaling curves de compute/data/performance & Extrapolar la fuerza impulsora a lo largo del tiempo y la escala \\
Revisión de límites & Examinar la métrica, la calidad de los datos y los puntos intermedios faltantes & Examinar las fuerzas ignoradas, los cuellos de botella físicos y el regime change \\
Paso a la acción & Diseñar el siguiente grupo de experimentos a escala & Agregar o retirar sesgos inductivos \\
\bottomrule
\end{tabularx}
\caption{El ciclo de investigación guiado por evidencia que comparten las dos conferencias.}
\end{table}

\subsection{Resumen de la sección}

Esta sección establece la forma de leer toda la clase. Jason pide al investigador que entrene su intuición sobre los datos; Hyung Won, que entrene su intuición histórica; ambos se oponen a saltar de un resultado aislado a conclusiones grandiosas. El contenido técnico que sigue repetirá una y otra vez la misma operación: descomponer el objetivo, entender las curvas, identificar los supuestos y explicar qué no puede concluirse a partir de la evidencia.
\section{Next-token Prediction: cómo un solo objetivo abarca una enorme cantidad de tareas}

La sección anterior presentó el método; esta sección aborda la primera intuición central de Jason. En apariencia, el entrenamiento de un modelo de lenguaje solo pide predecir el siguiente token, pero la probabilidad correcta de cada token puede depender de la sintaxis, el significado de las palabras, los hechos, las relaciones espaciales, la traducción, la aritmética e incluso el género del documento. Lo importante no es ponerle al preentrenamiento la etiqueta «multi-task», sino entender cómo el contexto forma de manera temporal los límites de cada tarea.

\subsection{De la distribución de probabilidad a la pérdida de entrenamiento}

Esta subsección completa primero el trasfondo matemático mínimo. Dados los tokens previos $x_{<t}=(x_1,\ldots,x_{t-1})$, el modelo produce la probabilidad condicional de cada token candidato del vocabulario; durante el entrenamiento solo se observa el token real $x_t$, y se aumenta su probabilidad. Al leer la slide, conviene distinguir primero entre «el modelo produce la distribución completa» y «los datos solo aportan un resultado observado».

\lecturefigure{jason-slide-004.jpg}{El modelo de lenguaje asigna a cada palabra candidata del vocabulario una probabilidad para la siguiente posición}{Jason Wei official deck, p.~4}

La log-verosimilitud negativa de un modelo de lenguaje autorregresivo puede escribirse como
\[
\mathcal{L}_{\mathrm{NLL}}(\theta)=-\frac{1}{T}\sum_{t=1}^{T}\log p_{\theta}(x_t\mid x_{<t}).
\]
Aquí $\theta$ representa los parámetros del modelo, $T$ es la longitud de la secuencia y $p_{\theta}(x_t\mid x_{<t})$ es la probabilidad que el modelo asigna al siguiente token real. Cuanto menor es la probabilidad, mayor es la penalización $-\log p$; el optimizador actualiza $\theta$ mediante el gradiente para que la continuation real de la distribución de entrenamiento sea más probable.

\begin{knowledgebox}{Qué es la perplexity}
La perplexity (perplejidad, PPL) es la exponencial de la entropía cruzada promedio:
\[
\operatorname{PPL}=\exp\!\left(\mathcal{L}_{\mathrm{NLL}}\right).
\]
Puede entenderse de manera aproximada como el «número efectivo de candidatos» que el modelo enfrenta en cada paso. Por ejemplo, una PPL de 10 no significa que el modelo elija realmente de manera uniforme entre diez palabras; es solo un resumen exponencial de la pérdida logarítmica, y tampoco equivale directamente a la exactitud factual, la tasa de aciertos en razonamiento ni las preferencias humanas.
\end{knowledgebox}

\subsection{El massive multi-task learning no es una metáfora}

Con la pérdida ya definida, el siguiente paso es preguntar qué determina la probabilidad de un token. Si el corpus es lo bastante grande, el modelo encontrará una y otra vez, en contextos distintos, sintaxis, semántica, hechos y formatos de tarea. Cada regularidad predecible aporta gradiente a la misma NLL; por ello, el objetivo unificado acumula en el espacio de parámetros muchas señales de aprendizaje locales.

\lecturefigure{jason-slide-005.jpg}{Intuición uno: la next-word prediction a gran escala es massive multi-task learning}{Jason Wei official deck, p.~5}

\begin{importantbox}{Definición de latent task}
Una latent task (tarea latente) es un subproblema de predicción que el contexto define de manera temporal, sin un task ID explícito. Por ejemplo, «The capital of Azerbaijan is» activa la recuperación de un hecho, «The movie was» puede activar una compleción de sentimiento y «The Spanish word for pretty is» activa una traducción. Comparten un mismo vocabulario y una misma pérdida, pero requieren cálculos intermedios distintos.
\end{importantbox}

\lecturefigure{jason-slide-006.jpg}{Las oraciones de preentrenamiento enseñan a la vez sintaxis, semántica, hechos, sentimiento, traducción, relaciones espaciales y aritmética}{Jason Wei official deck, p.~6}

\begin{knowledgebox}{Lectura de la figura: comparar primero «de dónde viene la supervisión»}
Ninguna fila etiqueta por separado el nombre de la tarea; la supervisión proviene de la continuation real. El ejemplo de sintaxis exige descartar palabras que no encajan sintácticamente; el de conocimiento del mundo exige colocar Baku por delante de London; el espacial exige seguir la posición de los personajes; el aritmético exige aprovechar las relaciones del enunciado. Lo que la figura respalda no es que «el modelo aprenda con seguridad cada capacidad», sino que «el corpus efectivamente aporta gradientes de probabilidad hacia estas capacidades».
\end{knowledgebox}

\teachervoice{Entre 00:04:59 y 00:07:03, Jason presenta varios ejemplos seguidos y al final enfatiza que puede haber millones de tareas así. El tono de la clase es importante: no afirma que los datos de preentrenamiento estén equilibrados de forma natural, sino que, mientras existan las oraciones pertinentes, el next-token objective puede extraer supervisión de ellas.}

\subsection{Las tareas pueden ser tan arbitrarias que resulta difícil nombrarlas}

Las tareas de la página anterior son todas muy «limpias» y fáciles de nombrar para una persona; esta subsección examina a propósito una continuation difícil de clasificar. Si una oración de Wikipedia pide predecir, en orden, un nombre de persona, una coma, un artículo y una palabra de profesión, cada posición recurre a una regularidad distinta. Las tareas no son un benchmark enumerado de antemano, sino preguntas momentáneas que plantea cada distribución condicional de la secuencia.

\lecturefigure{jason-slide-007.jpg}{Las tareas latentes de una misma oración pueden saltar de un hecho a la puntuación, y luego a regularidades de continuación difíciles de nombrar}{Jason Wei official deck, p.~7}

Si el token $t$-ésimo del corpus implica una tarea $z_t$, el objetivo de entrenamiento puede escribirse conceptualmente como
\[
\mathcal{L}=\mathbb{E}_{(x,z)\sim\mathcal{D}}\bigl[\ell(x_t,x_{<t};z_t)\bigr].
\]
Aquí $\mathcal{D}$ es la distribución de entrenamiento; $z_t$ no es una etiqueta que los datos proporcionen de forma explícita, sino una variable latente que introduce el analista para facilitar la comprensión, y $\ell$ es la pérdida de predicción en esa posición. Esta notación recuerda que un mismo token puede depender de varias regularidades a la vez; por eso, la descomposición en tareas es solo una herramienta explicativa y no garantiza que exista un límite único y correcto.

\begin{warningbox}{No hay que imaginar el pretraining corpus como un almacén de tareas ya depurado automáticamente}
El mismo mecanismo también aprende errores ortográficos, sesgos, correlaciones espurias, copias de plantillas y fuentes de baja calidad. En la sesión de Q\&A, Jason dice explícitamente que el entrenamiento no distingue de forma automática entre datos buenos y malos; los investigadores todavía deben revisar el origen de los datos, filtrar el contenido claramente poco confiable y comprobar, mediante evaluaciones independientes, si lo aprendido son regularidades transferibles o atajos de los datos.
\end{warningbox}

\begin{lstlisting}[caption={Pseudocódigo para descomponer un texto en latent-task probes}]
for position in document:
    context = tokens[:position]
    target = tokens[position]
    ask("What knowledge or rule raises target probability?")
    tag(grammar, fact, format, reasoning, noise, other)
    store(context, target, tags)
\end{lstlisting}

\subsection{Resumen de la sección}

La fuerza de la next-token prediction no se debe a que el nombre del objetivo sea complejo, sino a que el corpus reúne una gran cantidad de problemas de predicción condicional en un mismo objetivo. En cada paso, el modelo solo optimiza la probabilidad del token real, pero acumula en sus parámetros compartidos representaciones y patrones de cálculo que atraviesan distintas tareas. Esta perspectiva prepara la sección siguiente: si la loss global es una mezcla de las pérdidas de muchas tareas latentes, la curva global y las curvas de cada capacidad no tienen por qué tener la misma forma.
\section{Scaling Law y Emergence: cómo una pérdida total suave oculta cambios abruptos en las tareas}

La sección anterior descompuso el objetivo de entrenamiento en una mezcla de tareas latentes; esta sección estudia el eje de escala. La segunda y la tercera intuición de Jason no se contradicen: al aumentar la cantidad de cómputo, la cross-entropy total puede bajar de forma estable, mientras que la accuracy de un benchmark puede permanecer en cero durante mucho tiempo y luego subir de golpe. Para entender este fenómeno hay que observar al mismo tiempo la probabilidad continua, el peso de cada tarea y la regla de calificación discreta.

\subsection{Definición operativa mínima de scaling}

Este apartado explica primero el eje horizontal. En clase, el compute de entrenamiento se entiende de forma aproximada como el producto del tamaño del modelo por la cantidad de datos; el entrenamiento real también depende de la longitud de secuencia, del cómputo de las pasadas hacia adelante y hacia atrás, del aprovechamiento del hardware y de la receta de entrenamiento, pero esta simplificación basta para expresar la dirección de «invertir más cómputo efectivo». Lo importante es comparar una serie de puntos dentro del mismo paradigma de entrenamiento, no comparar de forma aislada los nombres de dos modelos.

\lecturefigure{jason-slide-008.jpg}{Intuición dos: el cómputo de entrenamiento que representan un modelo y unos datos más grandes reduce la loss de forma confiable}{Jason Wei official deck, p.~8}

Una aproximación común es
\[
C\propto N D,
\]
donde $C$ representa la cantidad de cómputo de entrenamiento, $N$ el número de parámetros sin contar los de embedding y $D$ el número de tokens de entrenamiento; la constante de proporcionalidad absorbe las operaciones hacia adelante y hacia atrás por token. Esta expresión no es un modelo de costo completo, pero le recuerda al investigador que aumentar solo los parámetros sin aumentar los datos, o aumentar solo los datos con un modelo de capacidad insuficiente, puede alejarlo de la región compute-optimal.

\lecturefigure{jason-slide-009.jpg}{La loss de un modelo de lenguaje sigue una tendencia predecible de ley de potencias respecto al cómputo}{Jason Wei official deck, p.~9}

El ajuste típico de una scaling law se escribe
\[
L(C)=L_{\infty}+A C^{-\alpha},\qquad A>0,\ \alpha>0.
\]
donde $L(C)$ es la pérdida de validación para un compute dado, $L_{\infty}$ es la cota inferior de la pérdida, irreducible o todavía no eliminada, $A$ controla la altura de la curva y $\alpha$ controla la pendiente de la mejora conforme crece el compute. En escala logarítmica, el intervalo de ley de potencias se aproxima a una recta, por lo que pueden usarse experimentos más pequeños para estimar la loss esperada de entrenamientos más grandes.

\begin{knowledgebox}{Lectura de la figura: confirmar primero los ejes, el rango y si hay saturación}
El eje horizontal abarca varios órdenes de magnitud y el eje vertical es la loss, donde menor es mejor. La evidencia visual decisiva no es que un punto vaya adelante, sino que la tendencia no se curva de forma apreciable hacia una horizontal en un rango muy amplio. Esto respalda que «dentro del intervalo observado, seguir aumentando el cómputo todavía rinde beneficios», pero no demuestra que la tendencia continúe para siempre ni que el beneficio alcance a compensar el costo de entrenamiento.
\end{knowledgebox}

\teachervoice{En 00:09:36--00:10:44, Jason recuerda que esta relación se llama ``law'' porque sigue siendo predecible a lo largo de unos siete órdenes de magnitud, no porque, como una ley física, nunca deje de cumplirse. Un informe de investigación debe dar a la vez el intervalo de ajuste y la incertidumbre de la extrapolación.}

\subsection{Por qué funciona el scaling: hay que separar la evidencia de las conjeturas}

El apartado anterior presentó una relación empírica; este discute los mecanismos posibles. Las dos explicaciones que se dan en clase se califican expresamente de hand-wavy: un modelo más grande puede memorizar más hechos de la cola larga y también puede representar funciones más complejas en una sola forward pass. Son hipótesis útiles, pero no pueden confirmarse solo a partir de la curva de loss.

\lecturefigure{jason-slide-010.jpg}{Intuición: el modelo pequeño memoriza de forma selectiva; el modelo grande da cabida a la cola larga y ejecuta cálculos más complejos}{Jason Wei official deck, p.~10}

\begin{warningbox}{Errores causales frecuentes al leer la figura}
«Un modelo más grande tiene menor loss» no basta para demostrar que «toda la reducción proviene de la memorization», ni que «la profundidad se encarga del razonamiento y el ancho de los hechos». El número de parámetros, la cobertura de los datos, la estabilidad de la optimización y el tiempo de entrenamiento cambian al mismo tiempo. Para poner a prueba un mecanismo hay que controlar variables, construir datos contrafactuales o medir directamente cómo cambian los hechos de distinta frecuencia y las tareas que exigen distinta profundidad de cálculo.
\end{warningbox}

\teachervoice{En 00:11:13--00:12:55, Jason dice dos veces que estas son solo hand-wavy answers. Conservar esta capa de incertidumbre importa más que ofrecer un mecanismo que suene completo: el hecho del curso es que la relación de scaling es robusta; «memoria de la cola larga + funciones más complejas» es una explicación candidata.}

\subsection{La loss total es una mezcla ponderada de las pérdidas por tarea}

Volvamos ahora a la perspectiva massive multi-task de la sección anterior. Si los distintos tokens del corpus requieren, respectivamente, gramática, hechos, sentimiento o matemáticas, la loss total se comporta como un promedio ponderado de las pérdidas de esas subdistribuciones. Las tareas fáciles pueden acercarse a la saturación mientras las difíciles todavía mejoran con rapidez; por eso, que la curva total sea suave no obliga a que cada componente lo sea.

\lecturefigure{jason-slide-011.jpg}{Intuición tres: la loss total baja de forma suave, pero una tarea posterior individual todavía puede mostrar un emergent pattern}{Jason Wei official deck, p.~11}

\lecturefigure{jason-slide-012.jpg}{Descomposición de la loss total en pérdidas por tarea, como grammar, world knowledge y math}{Jason Wei official deck, p.~12}

En términos conceptuales,
\[
L_{\mathrm{all}}(C)=\sum_{k=1}^{K} w_k L_k(C),
\qquad w_k\ge 0,\quad \sum_{k=1}^{K}w_k=1.
\]
donde $K$ es el número de familias de tareas latentes, $w_k$ es el peso de esa tarea entre los tokens de evaluación y $L_k(C)$ es la pérdida de la tarea cuando el compute es $C$. Aunque la pérdida $L_j$ de una tarea difícil baje rápidamente en un intervalo estrecho, si $w_j$ es muy pequeño la curva total todavía puede verse suave.

\begin{knowledgebox}{Qué es una emergent ability}
Esta clase adopta una definición conductual: un modelo pequeño obtiene en cierta tarea un resultado cercano al azar o de cero, mientras que un modelo más grande alcanza un resultado claramente mejor. Esta definición describe \textbf{la forma de la curva de evaluación}; no implica por sí misma que dentro del modelo haya surgido de repente un módulo completamente nuevo. Si la función de calificación tiene un umbral, usa exact match o exige acertar todos los pasos de un proceso de varios pasos, una mejora continua de la probabilidad también se traduce en un cambio abrupto aparente.
\end{knowledgebox}

\subsection{Las tareas fáciles se saturan y las difíciles arrancan tarde}

El apartado anterior dio la fórmula de la mezcla; este lee en la figura las curvas que la componen. La gramática puede ser ya suficientemente buena a una escala menor, de modo que seguir agrandando el modelo solo aporta mejoras mínimas; en cambio, las tareas difíciles, como las matemáticas, pueden cruzar el umbral de utilidad solo a una escala mayor. Al leer la figura hay que comparar las pendientes y la posición de la saturación, no interpretar cada línea como un módulo independiente.

\lecturefigure{jason-slide-013.jpg}{Formas de scaling distintas para la loss total, una tarea ya saturada y una tarea difícil}{Jason Wei official deck, p.~13}

Si la calificación de la tarea se obtiene a partir de una confianza continua $q(C)$ mediante un umbral $\tau$,
\[
S(C)=\mathbf{1}\{q(C)\ge \tau\},
\]
donde $S(C)$ es la calificación discreta, $q(C)$ es la confianza de éxito, que varía de forma suave con el compute, y $\tau$ es el umbral de calificación. Aunque $q(C)$ suba de forma suave, $S(C)$ salta de 0 a 1 al cruzar $\tau$. Esta es una de las fuentes de la emergence, pero no la única.

\lecturefigure{jason-slide-014.jpg}{Las tareas de BIG-Bench incluyen curvas smooth, flat, inverse, noisy y emergent-looking}{Jason Wei official deck, p.~14}

\begin{knowledgebox}{Lectura de la figura: la distribución de las tareas informa más que la dicotomía «emerge o no emerge»}
En la figura, cerca del 29\% de las tareas mejora de forma suave, el 22\% es casi plano dentro del rango observado, el 2\% muestra inverse scaling, el 13\% no tiene una relación clara con la escala y el 33\% es emergent-looking. Los porcentajes dependen del conjunto de modelos, de las tareas y de la metric; la conclusión más sólida es que las curvas por tarea adoptan muchas formas y que la tendencia única de la loss total no sustituye la medición tarea por tarea.
\end{knowledgebox}

\teachervoice{En 00:16:54--00:19:24, Jason enfatiza que, si solo se entrena hasta antes del punto de cambio abrupto, el investigador predecirá erróneamente que esa capacidad nunca aparecerá. En la sesión de Q\&A añade que la falta de puntos intermedios densos de modelos hace muy difícil determinar si ya había señal antes del cambio abrupto.}

\subsection{La emergence en el prompting y la salvedad sobre la metric}

Además de las curvas por tarea, el propio prompt también puede revelar diferencias de escala. Un modelo pequeño puede ser incapaz de extraer la regla de correspondencia a partir de unas cuantas demostraciones, mientras que uno grande puede aprovechar el contexto para resolver entradas nuevas; pero, tras observar un salto en exact-match, todavía hay que revisar la token-level probability, las metric continuas y la robustez frente al prompt, para no describir equivocadamente la geometría de la función de calificación como un mecanismo interno.

\lecturefigure{jason-slide-015.jpg}{Los ejemplos de entrada y salida en el prompt pueden activar un in-context behavior que depende de la escala}{Jason Wei official deck, p.~15}

\begin{warningbox}{«La emergence es un mirage» y «la capacidad es falsa» no son la misma afirmación}
En la sesión de Q\&A, Jason reconoce que cambiar la metric cambia la forma de la curva; aun así, sostiene que los comportamientos útiles que muestran los modelos grandes son reales. Una formulación rigurosa distingue tres niveles: si la distribución de salida del modelo cambia de forma continua, si la calificación de la evaluación cambia de golpe y si las tareas que el usuario puede resolver experimentan un cambio cualitativo. Los tres pueden cumplirse a la vez sin contradicción.
\end{warningbox}

\teachervoice{La respuesta de 00:29:53--00:30:21 no ofrece una etiqueta simple de bando; recomienda leer los artículos y juzgar por cuenta propia. Por eso estas notas no tratan la emergence como una cuestión de creencia, sino que consideran la metric sensitivity parte del diseño experimental.}

\subsection{Resumen de la sección}

La scaling law describe cómo la pérdida total mejora de forma estable con la cantidad de cómputo dentro del intervalo observado; la emergent-looking curve describe el comportamiento de una tarea específica después de pasar por la transformación de la metric. Ambas pueden surgir de la combinación de la mezcla de tareas, las diferencias de dificultad, la calificación por umbral y el muestreo disperso. El investigador debe informar a la vez la pérdida continua, la calificación por tarea, los puntos de tamaño de modelo y las barras de error, e indicar si cada explicación de un mecanismo es evidencia, hipótesis o intuición personal.
\section{Inverse y U-shaped Scaling: cómo se interfieren las subtareas ocultas}

La sección anterior trató la «aparición tardía»; esta sección trata el patrón «primero empeora y luego mejora». Este tipo de curva es la que con más facilidad se malinterpreta como una falla del entrenamiento; el ejemplo de Jason muestra que la tarea completa puede depender de la competencia entre varias capacidades. Un modelo mediano aprendió a corregir citas célebres, pero todavía no aprende a obedecer ``repeat after me'', por lo que obtiene peores resultados que un modelo más pequeño; un modelo más grande, que ya tiene ambas capacidades, vuelve a responder correctamente.

\subsection{Detrás de una curva puede haber varias estrategias}

Esta sección comienza por el fenómeno. Inverse scaling se refiere a que la puntuación de una tarea disminuye al aumentar la escala; U-shaped scaling, a que la puntuación primero baja y después vuelve a subir. La forma de la curva por sí sola no nos dice si se trata de una degradación de la optimización, de un conflicto entre los datos o de una competencia entre subtareas; por eso hay que construir un prompt interpretable que separe una por una las posibles estrategias.

\lecturefigure{jason-slide-016.jpg}{Intuición cuatro: elegir las tareas con cuidado produce inverse o U-shaped scaling}{Jason Wei official deck, p.~16}

\begin{knowledgebox}{Términos clave: inverse y U-shaped}
\begin{tabularx}{\textwidth}{L{0.23\textwidth} X}
\toprule
Término & Significado en el curso \\
\midrule
Inverse scaling & Al aumentar la escala del modelo o el cómputo, la metric de una tarea bien definida disminuye en lugar de aumentar. \\
U-shaped scaling & La misma metric primero baja y luego sube conforme crece la escala, lo que indica que distintas capacidades o estrategias dominan en distintos intervalos. \\
Task selection & Elegir prompts que permitan distinguir entre estrategias en competencia, en lugar de escoger solo el benchmark con el mayor desempeño promedio. \\
\bottomrule
\end{tabularx}
\end{knowledgebox}

\lecturefigure{jason-slide-017.jpg}{El modelo pequeño repite la palabra errónea, el mediano corrige la cita por su cuenta y el grande vuelve a seguir la instrucción}{Jason Wei official deck, p.~17}

Al leer la figura, primero hay que distinguir ``glib'' de ``gold''. El objetivo no es completar la cita célebre, sino obedecer la petición del usuario y repetir la oración escrita mal a propósito. El modelo pequeño la repite «por casualidad», porque no conoce la cita; el modelo mediano conoce la cita pero ignora la instruction; el modelo grande reconoce la cita y además es capaz de anteponer la instruction a su tendencia predeterminada de completar.

\subsection{Descomponer la tarea completa en repeat, repair e instruction following}

La página anterior presentó tres comportamientos; esta sección los formula explícitamente como subtareas. El modelo producirá ``glib'' de manera estable solo cuando sea capaz de repetir el texto, detectar el error en la cita y determinar la prioridad de la instrucción actual. Esto indica que la accuracy final es una función de una combinación de capacidades, y no una medición monótona de un único «grado de inteligencia».

\lecturefigure{jason-slide-018.jpg}{Las scaling curves de tres subtareas ocultas explican la forma de U del conjunto}{Jason Wei official deck, p.~18}

El comportamiento final puede simplificarse como
\[
\hat{y}=f\bigl(R(C),Q(C),I(C)\bigr),
\]
donde $R(C)$ representa la capacidad de repetir, $Q(C)$ la tendencia a corregir citas conocidas, $I(C)$ la capacidad de seguir instrucciones explícitas y $f$ la combinación de decisión de las tres. Un modelo pequeño puede tener $R$ alto, $Q$ bajo e $I$ bajo; un modelo mediano, $R$ alto, $Q$ alto e $I$ bajo; un modelo grande tiene las tres altas, y es $I$ la que determina que no corrija el texto del usuario.

\begin{warningbox}{No interpretes automáticamente la forma de U como «el modelo primero aprende mal y luego aprende bien»}
Una caída en la puntuación puede deberse a que una capacidad nueva se impone sobre una estrategia anterior, y no a una degradación general de las representaciones subyacentes. Un diagnóstico correcto requiere evaluar cada subcapacidad por separado y revisar qué objetivo del prompt debería tener prioridad. Una sola accuracy reduce a la misma etiqueta de correcto/incorrecto tanto «conoce la cita pero no obedece la instrucción» como «no conoce la cita en absoluto».
\end{warningbox}

\teachervoice{Entre 00:21:21 y 00:23:56, Jason dibuja una por una las curvas de las tres subtareas y luego regresa a la tarea completa para explicar la forma de U. Este es el procedimiento de análisis más reutilizable de toda la conferencia: ante una curva extraña, primero hay que preguntarse qué estrategias mezcla la metric.}

\subsection{Práctica de investigación: nunca grafiques un solo punto}

Una vez descompuestas las subtareas, el último paso es colocar también las intervenciones experimentales sobre el eje de escala. Que un fine-tuning, un filtrado de datos o un método de prompt supere a la baseline en un punto de presupuesto no demuestra que siga siendo eficaz al aumentar los datos. Se necesitan al menos varios puntos para determinar si la curva se satura, mantiene su pendiente o se separa cada vez más rápido; además, hay que repetir los experimentos con la misma definición de datos y el mismo protocolo de evaluación, para no confundir con una ventaja del método la fluctuación de un solo punto causada por la inicialización aleatoria o la selección de muestras.

\lecturefigure{jason-slide-019.jpg}{Recomendación general de investigación: graficar scaling curves en lugar de reportar solo la mejora en un punto}{Jason Wei official deck, p.~19}

\begin{lstlisting}[caption={Comparación de las scaling curves de una intervención de investigación y de la baseline}]
budgets = [0.25, 0.5, 1.0, 2.0]
for budget in budgets:
    baseline = train(config="baseline", budget=budget)
    method = train(config="intervention", budget=budget)
    log(budget, evaluate(baseline), evaluate(method))
fit_curve_and_report_uncertainty()
\end{lstlisting}

\begin{importantbox}{Una scaling curve responde al menos tres preguntas}
\begin{enumerate}
  \item \textbf{Data efficiency}: ¿el método solo adelanta el mismo desempeño a un presupuesto menor?
  \item \textbf{Asymptote}: ¿la ventaja se satura pronto, de modo que invertir más ya no la amplía?
  \item \textbf{Slope}: ¿el método cambia la pendiente de crecimiento, de modo que el beneficio es mayor a mayor escala?
\end{enumerate}
Un resultado de un solo punto no permite distinguir entre estos tres casos.
\end{importantbox}

\teachervoice{La conclusión de 00:23:59 a 00:25:51 es muy sencilla: ``just plot scaling curves''. No se pide que cada estudiante entrene un frontier model, sino que haga experimentos escalonados dentro del rango de presupuesto que pueda costear, para no tomar un punto fortuito por una tendencia.}

\subsection{Resumen de la sección}

Inverse y U-shaped scaling suelen indicar que detrás de la metric hay estrategias en competencia. Mediante prompts interpretables, evaluaciones de subcapacidades y curvas con varios presupuestos, los investigadores pueden reformular un «fenómeno extraño» como una hipótesis verificable. Aquí se cierra el ciclo completo de la primera conferencia de Jason: observar los datos, identificar las tareas latentes, descomponer la loss, graficar scaling curves y señalar con claridad la metric y los límites de la evidencia.
\section{Compute, Bitter Lesson y el ciclo de vida del sesgo inductivo}

Ahora entramos en la línea principal de Hyung Won. En la primera mitad las curvas se trazan en función de la escala del modelo; en la segunda, en función del tiempo histórico. El eje horizontal común sigue siendo el cómputo disponible y la capacidad de los algoritmos. El juicio central de Hyung Won es que la caída sostenida del costo del cómputo, junto con el associated scaling, es la fuerza dominante de la investigación en IA. Por eso, al diseñar una estructura no basta con ver si hoy es más precisa; también hay que ver si limitará seguir ampliando la escala mañana.

\subsection{Por qué el compute per dollar se toma como fuerza dominante}

Esta sección empieza por la evidencia. El eje vertical de la figura no es el número de transistores, sino la capacidad de cómputo que puede comprarse con una cantidad fija de dinero; el eje horizontal abarca más de un siglo. A partir de esta tendencia de largo plazo, Hyung Won propone una estrategia de investigación: no competir contra una tendencia externa fuerte y persistente, sino diseñar métodos que la aprovechen. Al mismo tiempo, hay que incluir el hardware, los algoritmos, la precisión numérica y la eficiencia del software en el «cómputo efectivo»; de lo contrario, se confunde la métrica de un solo componente con la trayectoria completa de la investigación.

\lecturefigure{hyung-slide-007.jpg}{La capacidad de cómputo que se obtiene a costo fijo crece a largo plazo, aproximadamente diez veces cada cinco años}{Hyung Won Chung official deck, p.~7}

Si la capacidad de cómputo crece aproximadamente diez veces cada cinco años, puede escribirse
\[
C(t)=C_0\,10^{t/5},
\]
donde $C(t)$ es la cantidad de cómputo que puede comprarse con un presupuesto fijo $t$ años después del punto de partida, y $C_0$ es la cantidad de cómputo en el punto de partida. Esta expresión solo aproxima la figura de la clase y no garantiza que el mismo exponente se mantenga en el futuro; la energía, la fabricación, la interconexión, el capital y los algoritmos pueden modificar la tasa de crecimiento efectiva.

\begin{knowledgebox}{Lectura de la figura: compute availability no es lo mismo que Moore's law}
En la sesión de preguntas y respuestas, el ponente distingue explícitamente entre la densidad de transistores y el cómputo disponible. Las GPU, la baja precisión, los aceleradores especializados, la pila de software y el codiseño con el hardware pueden seguir mejorando el rendimiento del entrenamiento, aunque el escalamiento tradicional de los transistores se desacelere. La pregunta pertinente es «cuánto cómputo de entrenamiento efectivo puede realizarse con un costo dado», no la métrica de un solo componente.
\end{knowledgebox}

\teachervoice{Entre 00:38:24 y 00:39:27 la formulación es «no compitas con esta tendencia; aprovéchala tanto como puedas». Es una heurística para elegir líneas de investigación, no una garantía de crecimiento exponencial ilimitado; después de 01:15:17, en la sesión de preguntas y respuestas, el ponente menciona por iniciativa propia que la energía y los límites físicos podrían convertirse en nuevos cuellos de botella.}

\subsection{Bitter Lesson: imponerle menos a la máquina la forma de pensar que creemos entender}

La sección anterior mostró que la escala es una tendencia fuerte; esta explica por qué el exceso de estructura diseñada a mano entra en conflicto con ella. Los investigadores suelen codificar en el modelo su propia explicación del pensamiento humano; con poco cómputo, estas estructuras pueden parecer atajos, pero limitan la exploración de otras representaciones por parte del modelo. La versión de la Bitter Lesson que se presenta en clase es esta: el progreso a largo plazo ha surgido una y otra vez de métodos más generales, con supuestos más débiles, sumados a más datos y más cómputo.

\lecturefigure{hyung-slide-008.jpg}{«Enseñar a la máquina a pensar como creemos que pensamos» puede imponer una estructura equivocada}{Hyung Won Chung official deck, p.~8}

\lecturefigure{hyung-slide-009.jpg}{Bitter Lesson: métodos más generales, supuestos de modelado más débiles y mayor escala}{Hyung Won Chung official deck, p.~9}

\begin{importantbox}{Qué es el inductive bias}
El inductive bias (sesgo inductivo) es el supuesto estructural por el cual, con datos limitados, un modelo prefiere cierta clase de funciones, representaciones o formas de interacción. La localidad de la convolución, la separación de parámetros entre encoder y decoder, la attention bidireccional y ciertas invariancias específicas son sesgos. Un sesgo puede mejorar la eficiencia muestral, pero también excluye parte de las soluciones; su valor depende de los datos, el cómputo, la tarea y las restricciones de despliegue del momento.
\end{importantbox}

\teachervoice{Entre 00:39:41 y 00:41:43, Hyung Won dice que los investigadores suelen intentar modelar «cómo creemos que pensamos», sin entender realmente los mecanismos cognitivos de bajo nivel. La postura de la clase es deliberadamente tajante; estas notas la conservan como una advertencia de diseño, no como una demostración de que toda estructura previa sea errónea.}

\subsection{Cuanta menos estructura, más scalable, pero no necesariamente mejor hoy}

Si se compara según la cantidad de estructura, los métodos con más estructura suelen partir de un nivel más alto en el intervalo de poco cómputo, pero pueden saturarse antes; los métodos con menos estructura son difíciles de entrenar al principio, pero dan más libertad al modelo y, al crecer la escala, pueden superar a los primeros. Aquí ``scalable'' significa que la curva de crecimiento puede seguir mejorando, no solo que el método pueda ejecutarse en más GPU.

\lecturefigure{hyung-slide-010.jpg}{Más estructura mejora el desempeño con poco cómputo, pero puede reducir la escalabilidad a largo plazo}{Hyung Won Chung official deck, p.~10}

\begin{knowledgebox}{Lectura de la figura: comparar el punto de cruce en lugar de calificar la estructura como buena o mala}
A la izquierda del presupuesto actual, el método azul puede ser claramente mejor; el método rojo solo lo supera con más compute. La decisión de investigación depende de dónde esté el punto de despliegue, de si puede asumirse el costo de exploración y de si en el futuro realmente se llegará al punto de cruce. Lo que la figura respalda es una regime-dependent choice, no «elegir siempre la menor estructura posible».
\end{knowledgebox}

\lecturefigure{hyung-slide-013.jpg}{En el compute regime actual conviene elegir un sesgo que funcione, sin olvidar retirarlo en el futuro}{Hyung Won Chung official deck, p.~13}

\lecturefigure{hyung-slide-014.jpg}{Dados el cómputo, los datos, los algoritmos y la arquitectura, existe un inductive bias óptimo de manera temporal}{Hyung Won Chung official deck, p.~14}

Si la intensidad de la estructura se denota por $b$ y el estado actual de los recursos por $r$, el desempeño se escribe como
\[
b^{\star}(r)=\arg\max_b\; P(b;r).
\]
Aquí $P(b;r)$ es el desempeño del sesgo $b$ en el estado de recursos $r$, y $b^{\star}(r)$ es la estructura óptima actual. A medida que cambian el compute, los data y los algorithm que contiene $r$, $b^{\star}$ también se desplaza; por eso, una decisión de arquitectura que alguna vez fue correcta no debe congelarse para siempre.

\teachervoice{La salvedad central entre 00:42:28 y 00:43:51 es que no se puede esperar desde el principio al método más general, porque podría no funcionar en absoluto con el presupuesto actual. Lo correcto es «incorporar ahora el sesgo útil y dejar registrado explícitamente cuándo revisarlo y retirarlo en el futuro».}

\subsection{Por qué los incentivos de la investigación favorecen agregar estructura y no eliminarla}

El ciclo de vida de la estructura también tiene causas sociales. Un módulo, un término de pérdida o un mecanismo nuevo se convierte fácilmente en la contribución de un artículo; demostrar que un módulo complejo puede eliminarse a mayor escala suele parecer poco novedoso. Así, la comunidad acumula supuestos diseñados para un regime anterior, pero carece de una práctica sistemática de eliminación. Esta sección considera los incentivos de la investigación como parte de la evolución de las arquitecturas, porque qué experimentos reciben recursos y qué resultados negativos se publican influye directamente en cuánto tiempo se conservan los sesgos antiguos.

\lecturefigure{hyung-slide-015.jpg}{El método que es mejor a largo plazo suele ser peor en el presente}{Hyung Won Chung official deck, p.~15}

\begin{warningbox}{Una ventaja de corto plazo en el leaderboard no equivale a valor de investigación a largo plazo}
Si un método solo lleva la delantera con modelos pequeños, pocos datos o un benchmark específico, hay que seguir preguntando: ¿la ventaja proviene de un prior más fuerte? ¿Se satura al aumentar la escala? ¿La complejidad de implementación dificulta entrenamientos más grandes? Estas preguntas no pueden responderse con la clasificación de un único punto actual.
\end{warningbox}

\teachervoice{Entre 00:43:51 y 00:44:47, Hyung Won afirma sin rodeos que los incentivos de publicación fomentan agregar estructura, pero rara vez recompensan eliminarla. También dice que el método que es mejor a largo plazo casi inevitablemente parece peor hoy, porque un sistema de aprendizaje con más libertad es más caótico en el intervalo de pocos recursos.}

\lecturefigure{hyung-slide-016.jpg}{Resumen de la etapa: la fuerza dominante es el cómputo más barato y el associated scaling}{Hyung Won Chung official deck, p.~16}

Esta diapositiva de resumen no niega todas las demás fuerzas; anuncia el sistema de coordenadas del análisis siguiente: explicar, una por una, las estructuras adicionales de la arquitectura Transformer como inductive bias propios de ciertas condiciones históricas, y preguntar si la escala actual todavía las necesita. Así, el diagrama de arquitectura deja de ser una clasificación estática y se convierte en un objeto sobre el que se puede razonar a lo largo de la scale trajectory.

\subsection{Resumen de la sección}

Esta sección estableció el marco de juicio de Hyung Won: la caída sostenida del costo del cómputo da a los métodos generales un gran potencial a largo plazo; el sesgo inductivo mejora la eficiencia cuando hay pocos recursos, pero puede limitar el escalamiento posterior. La conclusión correcta no es «cuanta más estructura, mejor» ni «cuanta menos estructura, mejor», sino estimar continuamente el sesgo óptimo del regime actual y dejar abierto el camino para futuros experimentos de eliminación.
\section{Tres tipos de Transformer: a partir de la base común de los sequence models}

La sección anterior explicó por qué conviene estudiar la estructura; esta sección primero deja clara la estructura. Encoder-decoder, encoder-only y decoder-only no son tres invenciones sin relación entre sí: las tres convierten una entrada discreta en una secuencia de vectores y después usan attention para modelar las interacciones entre los elementos. Las diferencias provienen sobre todo de si los parámetros se comparten, de cómo la attention mask limita el alcance visible y de cómo se entrena y se usa la salida.

\subsection{La interfaz de tarea de las tres arquitecturas}

Esta sección comienza con el panorama general y después entra paso a paso en las figuras. Encoder-decoder asigna dos stacks, uno para la entrada y otro para la salida; encoder-only comprime la entrada en una representación y por lo común necesita un task-specific head; decoder-only concatena la entrada y el objetivo y genera token por token en un solo causal stack. Cada interfaz supone implícitamente un tipo distinto de tarea.

\lecturefigure{hyung-slide-017.jpg}{Las tres arquitecturas del Transformer: encoder-decoder, encoder-only y decoder-only}{Hyung Won Chung official deck, p.~17}

\begin{knowledgebox}{Glosario: qué problema resuelve cada arquitectura}
\begin{tabularx}{\textwidth}{L{0.22\textwidth} X}
\toprule
Arquitectura & Mecanismo central y relación con el curso \\
\midrule
Encoder-decoder & El encoder codifica la entrada de forma bidireccional; el decoder genera el objetivo mediante causal self-attention y cross-attention. Es adecuada para sequence-to-sequence explícito. \\
Encoder-only & Interacción bidireccional sobre toda la secuencia; la salida suele usarse para clasificación, recuperación o token labeling. Para generar una secuencia se necesita un mecanismo de generación adicional. \\
Decoder-only & La entrada y las salidas previas están en la misma secuencia causal; comparte parámetros y tiene una interfaz unificada, por lo que se adapta de forma natural a la autoregressive generation y a la continuación en varios turnos. \\
\bottomrule
\end{tabularx}
\end{knowledgebox}

\subsection{De caracteres a una secuencia de vectores, y de ahí a la interaction}

Antes de clasificar las arquitecturas hay que entender el pipeline común. Primero, el tokenizer convierte el texto en token IDs; después, la embedding table asigna a cada ID un vector de dimensión $d_{\text{model}}$; por último, el sequence model hace que las distintas posiciones intercambien información según el attention pattern. El modelo procesa una secuencia de vectores, no directamente el «significado de las palabras».

\lecturefigure{hyung-slide-018.jpg}{La secuencia de caracteres sin procesar es el punto de partida del pipeline del Transformer}{Hyung Won Chung official deck, p.~18}

\lecturefigure{hyung-slide-021.jpg}{Pipeline completo: tokenización, embedding y sequence model}{Hyung Won Chung official deck, p.~21}

Si el token ID es $i_t$ y la embedding matrix es $E\in\mathbb{R}^{V\times d}$, entonces
\[
h_t^{(0)}=E[i_t]+p_t.
\]
Aquí $V$ es el tamaño del vocabulario, $d$ es la hidden dimension, $E[i_t]$ es el vector del token número $i_t$, $p_t$ es la representación de posición y $h_t^{(0)}$ es el estado que entra en la primera capa. La tokenización determina la longitud de la secuencia y los límites de segmentación; el embedding determina la representación continua inicial. Ambos influyen en el costo del cómputo posterior.

\teachervoice{Entre 00:46:35 y 00:47:50, Hyung Won primero define el Transformer como un sequence model: los elementos de entrada pueden ser palabras, imágenes u otros objetos, y lo central es modelar la interaction entre los elementos. Este orden evita separar la fórmula de attention de la interfaz de datos.}

\subsection{Encoder-decoder: tres papeles de la attention}

Pasamos ahora al Transformer original. Al leer la figura, primero hay que separar los tres flujos de información: la self-attention bidireccional dentro del encoder, la self-attention causal dentro del decoder y la cross-attention con la que el decoder consulta la representación del encoder. Cada arista corresponde a un permiso de «qué posición puede leer qué posición».

\lecturefigure{hyung-slide-022.jpg}{Arquitectura encoder-decoder completa y los tres papeles de la attention}{Hyung Won Chung official deck, p.~22}

\begin{importantbox}{Self-attention y cross-attention}
Dados una query $Q$, una key $K$ y un value $V$, la scaled dot-product attention es
\[
\operatorname{Attn}(Q,K,V)=\operatorname{softmax}\!\left(\frac{QK^{\top}}{\sqrt{d_k}}+M\right)V.
\]
Aquí $d_k$ es la dimensión de la key y $M$ es la máscara de visibilidad. En la self-attention, $Q,K,V$ provienen de la misma secuencia; en la cross-attention, $Q$ proviene del decoder y $K,V$ provienen de la salida del encoder.
\end{importantbox}

\lecturefigure{hyung-slide-023.jpg}{El encoder usa bidirectional self-attention para modelar la entrada completa}{Hyung Won Chung official deck, p.~23}

La bidirectional attention (atención bidireccional) permite que cualquier posición de la entrada lea cualquier otra posición. Si la longitud de la secuencia es $n$, la máscara de visibilidad puede escribirse como $M_{ij}=0$; no hay ocultamiento del futuro. Es adecuada para una entrada que se da completa de una vez, porque todos los tokens se conocen al codificar; pero también implica que un cambio al final de la entrada afecta la representación de todas las posiciones anteriores.

\lecturefigure{hyung-slide-025.jpg}{El decoder usa causal self-attention y solo lee la posición actual y el historial a su izquierda}{Hyung Won Chung official deck, p.~25}

La causal attention (atención causal) usa una máscara triangular superior para impedir la lectura del futuro:
\[
M_{ij}=\begin{cases}
0,&j\le i,\\
-\infty,&j>i.
\end{cases}
\]
Aquí $i$ es la posición de la query actual y $j$ es la posición de la key que se lee. Después del softmax, el peso de las posiciones futuras se vuelve cero; por eso, durante el entrenamiento todas las posiciones pueden calcularse en paralelo y, durante la inferencia, se mantiene la condición autoregressive.

\lecturefigure{hyung-slide-027.jpg}{La cross-attention conecta las posiciones del objetivo en el decoder con la representación de la entrada en el encoder}{Hyung Won Chung official deck, p.~27}

\lecturefigure{hyung-slide-028.jpg}{Construcción final del encoder-decoder: la entrada se codifica de forma bidireccional y el objetivo se genera de forma causal}{Hyung Won Chung official deck, p.~28}

\begin{knowledgebox}{Lectura de la figura: cómo fluye una ruta de traducción automática}
La entrada ``That is good'' primero interactúa de forma bidireccional dentro del encoder y forma una representación contextualizada de cada token de origen. Al generar ``Das ist gut'', el decoder solo puede leer el prefijo del objetivo ya generado y, al mismo tiempo, accede a la última capa del encoder mediante cross-attention. La figura respalda una estructura de permisos de información; no indica que los pesos de attention equivalgan por naturaleza a una explicación como alineación.
\end{knowledgebox}

\subsection{Diferencias clave entre encoder-only y decoder-only}

La sección anterior usó dos stacks; esta sección examina los dos extremos. Encoder-only produce una representación bidireccional global; para clasificar, suele leer un token especial y añadir una task-specific layer. Decoder-only, en cambio, unifica la entrada y la salida en una sola secuencia y usa el mismo conjunto de parámetros y causal self-attention tanto para comprender el contexto como para generar.

\lecturefigure{hyung-slide-032.jpg}{Forma completa de encoder-only: representación bidireccional, tokens especiales y cabeza de tarea}{Hyung Won Chung official deck, p.~32}

\begin{warningbox}{Encoder-only no significa «no puede producir ninguna secuencia»}
Esta diapositiva se refiere a la interfaz por omisión: un solo encoder stack no trae integrada una autoregressive generation path. Encima de él puede agregarse un decoder, un span head, un rellenado iterativo de huecos u otro mecanismo de generación; el costo es que el sistema deja de ser la estructura encoder-only mínima de la figura. El curso compara los supuestos integrados, no la capacidad absoluta en un sentido lógico.
\end{warningbox}

\lecturefigure{hyung-slide-037.jpg}{Forma completa de decoder-only: entrada y objetivo concatenados, parámetros compartidos y causal attention}{Hyung Won Chung official deck, p.~37}

La interfaz unificada de decoder-only puede escribirse como
\[
p(y\mid x)=\prod_{t=1}^{|y|}p(y_t\mid x,y_{<t}).
\]
Aquí $x$ es la secuencia de tokens de entrada, $y$ es la secuencia objetivo y $y_{<t}$ es el prefijo del objetivo ya generado. La entrada y el objetivo usan el mismo stack, pero el loss por lo común puede calcularse solo en las posiciones del objetivo. «Compartir parámetros» no significa que el modelo no pueda distinguir los papeles, porque la posición, los tokens separadores y el propio contexto siguen aportando condiciones.

\subsection{Resumen de la sección}

Los tres tipos de Transformer comparten la base de tokenización, embedding, attention y MLP; las diferencias principales son la visibilidad de la información, la agrupación de los parámetros y la interfaz de tarea. Una vez comprendidas estas diferencias, la siguiente sección ya no trata encoder-decoder y decoder-only como categorías inconmensurables, sino que transforma de manera continua la primera en la segunda mediante cuatro operaciones, con lo que se pone al descubierto el verdadero sesgo inductivo.
\section{Transformación en cuatro pasos: reescribir de forma continua un Encoder-decoder como Decoder-only}

La sección anterior presentó los diagramas estáticos de las arquitecturas; esta sección desarrolla la derivación central de Hyung Won. El punto de partida es un encoder-decoder estándar y el de llegada, un decoder-only; cada paso elimina una sola diferencia y actualiza la tabla comparativa. El valor didáctico de la derivación no está en recomendar una implementación concreta, sino en reformular la afirmación «las arquitecturas son distintas» como cuatro hipótesis que pueden probarse por separado.

\subsection{Punto de partida: cuatro tipos de diferencias}

Esta subsección fija primero las dimensiones de comparación: la cross-attention adicional, si los parámetros del encoder y del decoder están separados, en qué capa se conecta target-to-input y si la self-attention de la entrada es bidireccional. Cada paso posterior modifica solo uno de estos elementos, para no mezclar varios cambios en una misma ablation.

\lecturefigure{hyung-slide-040.jpg}{Punto de partida de la transformación: comparación lado a lado de un encoder-decoder estándar y un decoder-only}{Hyung Won Chung official deck, p.~40}

\begin{importantbox}{Cuatro hipótesis comprobables}
\begin{enumerate}
  \item ¿La entrada y el objetivo necesitan attention roles independientes?
  \item ¿La entrada y el objetivo necesitan parámetros independientes?
  \item ¿Cada capa del decoder solo puede leer la capa final del encoder?
  \item ¿La entrada debe usar necesariamente bidirectional attention?
\end{enumerate}
Solo al separarlas puede saberse de qué elemento de la estructura proviene una diferencia de desempeño.
\end{importantbox}

\subsection{Paso uno: compartir los parámetros de cross-attention y self-attention}

Tanto la cross-attention como la self-attention usan query, key, value y output projection; la diferencia principal es de dónde provienen los tensores de entrada. Por ello, el primer paso no elimina ninguna conexión, sino que hace que ambos roles compartan los parámetros de proyección, de modo que un solo attention module se encargue a la vez de la interacción dentro de la secuencia y entre secuencias.

\lecturefigure{hyung-slide-041.jpg}{Paso uno: compartir los parámetros de cross-attention y self-attention}{Hyung Won Chung official deck, p.~41}

Si originalmente se tienen $W_Q^{\mathrm{self}},W_K^{\mathrm{self}},W_V^{\mathrm{self}}$ y $W_Q^{\mathrm{cross}},W_K^{\mathrm{cross}},W_V^{\mathrm{cross}}$, compartir parámetros equivale a imponer la restricción
\[
W_{\bullet}^{\mathrm{self}}=W_{\bullet}^{\mathrm{cross}},
\quad \bullet\in\{Q,K,V,O\}.
\]
Aquí $Q,K,V,O$ representan, respectivamente, query, key, value y output projection. Las fuentes de las conexiones pueden seguir siendo distintas, pero el modelo ya no puede conservar transformaciones lineales completamente independientes para los dos roles.

\lecturefigure{hyung-slide-042.jpg}{Comparación tras el paso uno: la cross-attention independiente se convierte en una self-attention que cumple ambos roles}{Hyung Won Chung official deck, p.~42}

\begin{knowledgebox}{Lectura de la figura: compartir parámetros no equivale a compartir entradas}
Lo que cambia en la tabla es la module identity, no la visibilidad. El decoder todavía puede aplicar masks distintas a su propio historial y a las representaciones del encoder; lo único que coincide son las matrices de proyección. Si en un experimento se cambian a la vez la mask y los parámetros, es difícil determinar si la mejora proviene de compartir representaciones o del patrón de conexión.
\end{knowledgebox}

\subsection{Paso dos: compartir los parámetros del encoder y del decoder}

El primer paso unifica los attention roles; el segundo unifica los dos stacks. Si tanto la entrada como el objetivo son simplemente token sequences, usar parámetros completamente independientes equivale a suponer que ambos tipos de secuencia requieren espacios de representación distintos. Compartir parámetros, en cambio, exige que un mismo conjunto de layers procese ambas, y deja que el contexto y la mask expresen la diferencia. Por eso esta subsección no solo compara el número de parámetros, sino también dos decisiones de modelado: si la diferencia de roles debe escribirse en los pesos o expresarse dinámicamente según la secuencia.

\lecturefigure{hyung-slide-044.jpg}{Paso dos: compartir los parámetros de las layers del encoder y del decoder}{Hyung Won Chung official deck, p.~44}

\lecturefigure{hyung-slide-045.jpg}{Comparación tras el paso dos: los separate parameters se convierten en shared parameters}{Hyung Won Chung official deck, p.~45}

\begin{warningbox}{Compartir parámetros es una suposición a priori más débil sobre la tarea, y también puede causar interferencia}
Compartir reduce los parámetros y favorece la reutilización de conocimiento entre idiomas y entre roles, pero cuando las distribuciones de la entrada y del objetivo difieren mucho, también puede generar competencia por la capacidad. La pregunta correcta no es si «compartir siempre es más avanzado», sino si la tarea actual todavía requiere especialización y si el beneficio de esa especialización se mantiene al aumentar la capacidad total.
\end{warningbox}

\teachervoice{La explicación de 00:55:17--00:55:41 es muy directa: un decoder-only tiene un solo stack, así que comparte parámetros de forma natural. Hyung Won usa este contraste para cuestionar si la separación del encoder-decoder sigue siendo necesaria, no para afirmar que compartir no tenga costo.}

\subsection{Paso tres: alinear las representaciones de la misma capa del decoder y del encoder}

La cross-attention estándar suele hacer que cada decoder layer lea la capa final del encoder. En cambio, la self-attention de misma capa del decoder-only hace que la capa $\ell$ lea las representaciones de los tokens del prefijo en esa misma capa. El tercer paso cambia la conexión a una layer-aligned, de modo que la granularidad de las representaciones se acerque a la de un solo stack.

\lecturefigure{hyung-slide-047.jpg}{Paso tres: decoder layer $\ell$ attends to encoder layer $\ell$}{Hyung Won Chung official deck, p.~47}

\lecturefigure{hyung-slide-048.jpg}{Comparación tras el paso tres: la final-layer cross-attention se convierte en una conexión entre capas del mismo nivel}{Hyung Won Chung official deck, p.~48}

\begin{knowledgebox}{Lectura de la figura: por qué el layer index puede representar la granularity}
Las redes profundas suelen presentar representaciones jerárquicas: las capas inferiores tienden a capturar la forma local y las superiores, la semántica abstracta. Si la primera capa del decoder solo puede leer la capa final del encoder, tiene que recuperar información de bajo nivel a partir de una representación muy abstracta; la conexión entre capas del mismo nivel permite que cada etapa lea entradas de granularidad similar. Esta explicación es una intuición de diseño; si realmente se forma un cuello de botella aún debe comprobarse con ablations a distintas profundidades.
\end{knowledgebox}

\subsection{Paso cuatro: cambiar la self-attention del encoder a causal}

Los tres primeros pasos ya unificaron el module, los parámetros y las conexiones entre capas; la principal diferencia restante es si la secuencia de entrada tiene visibilidad bidireccional interna. Al cambiar la mask del encoder a causal, los tokens anteriores ya no leen los posteriores; si se concatenan la entrada y el objetivo, el grafo de cómputo de los dos stacks se aproxima al de un único stack decoder-only.

\lecturefigure{hyung-slide-050.jpg}{Paso cuatro: cambiar la bidirectional self-attention del encoder a causal}{Hyung Won Chung official deck, p.~50}

\lecturefigure{hyung-slide-051.jpg}{Tabla de diferencias tras completar los cuatro pasos: entre ambas arquitecturas solo quedan pequeñas diferencias de implementación}{Hyung Won Chung official deck, p.~51}

\begin{lstlisting}[caption={Pseudocódigo conceptual de la architecture transformation en cuatro pasos}]
model = encoder_decoder()
share(model.cross_attention, model.self_attention)
share(model.encoder_layers, model.decoder_layers)
connect_decoder_layer_to_matching_encoder_layer(model)
set_encoder_mask(model, mode="causal")
concatenate_input_and_target(model)
\end{lstlisting}

\teachervoice{En 00:56:45--00:57:09, Hyung Won sostiene que, tras completar los cuatro pasos, la diferencia entre ambas arquitecturas con la misma tarea y los mismos datos podría ser menor que el ruido de repetir el entrenamiento. Se trata de un juicio empírico, no de un teorema formal presentado en el curso; una verificación real aún requiere controlar el número de parámetros, la inicialización, el presupuesto de entrenamiento y el orden de los datos.}

\subsection{Resumen de la sección}

La transformación en cuatro pasos reduce las etiquetas de arquitectura a cuatro hipótesis locales: attention role, separación de parámetros, conexiones entre capas y direccionalidad de la entrada. Todas pueden probarse por separado, y el valor de cada una puede cambiar con la escala. La siguiente sección explicará, una por una, por qué estas estructuras eran razonables en la traducción automática y en las primeras tareas, y por qué deben reconsiderarse en los modelos de lenguaje de propósito general, el diálogo de varios turnos y las redes más profundas.
\section{Por qué la estructura funcionó en su momento: el experimento natural de la traducción y el Instruction Tuning}

La sección anterior completó la transformación formal; esta sección se pregunta por la motivación histórica. La estructura adicional del encoder-decoder no es una complicación arbitraria: en 2017 la tarea central era la traducción automática, en la que la entrada y el objetivo provienen de idiomas distintos; más tarde, los datos académicos de instruction-tuning también solían presentar entradas largas y respuestas cortas. Cuando la estructura coincide con la distribución de los datos, el sesgo inductivo aumenta de manera notable la eficiencia con un presupuesto limitado.

\subsection{Tres estructuras adicionales}

Esta sección revisa primero la lista completa y después se concentra en el primer elemento. En comparación con decoder-only, encoder-decoder supone que la entrada y el objetivo son lo bastante distintos, que el objetivo debe leer una entrada ya codificada por completo y que los token de la entrada se prestan a una interacción totalmente bidireccional. Los tres supuestos corresponden, respectivamente, a los parámetros, a las conexiones entre capas y a la attention mask.

\lecturefigure{hyung-slide-052.jpg}{Las tres estructuras adicionales de encoder-decoder en comparación con decoder-only}{Hyung Won Chung official deck, p.~52}

\lecturefigure{hyung-slide-053.jpg}{Primer supuesto: la entrada y el objetivo son tan distintos que justifican parámetros independientes}{Hyung Won Chung official deck, p.~53}

\begin{importantbox}{Tres preguntas para juzgar si un inductive bias sigue siendo válido}
\begin{enumerate}
  \item ¿Con qué distribución de datos o restricción de la tarea coincidía originalmente?
  \item ¿Los datos, objetivos y formas de interacción actuales siguen cumpliendo esa restricción?
  \item Si se retira la estructura y se aumenta la escala, ¿en qué punto del presupuesto se cruzan las curvas de rendimiento?
\end{enumerate}
Responder solo la primera pregunta confunde la validez histórica con una corrección permanente.
\end{importantbox}

\subsection{Traducción automática: los parámetros independientes eran naturales}

En 2017, el Transformer estaba orientado a la source-to-target translation. La entrada es una oración en el idioma de origen y el objetivo es una salida en otro idioma; ambos difieren en vocabulario, orden de las palabras y papel en la generación. En un entorno que solo optimiza la traducción, dejar que el encoder se especialice en el idioma de origen y el decoder en el idioma de destino es un supuesto previo eficiente. Esta sección reconoce primero esa validez histórica y después compara qué condiciones cambió el preentrenamiento general, para no deducir, a partir de las interfaces de tareas actuales, que el diseño antiguo «era erróneo desde el principio».

\lecturefigure{hyung-slide-054.jpg}{La diferencia de distribución entre source y target en la traducción automática justifica los parámetros independientes}{Hyung Won Chung official deck, p.~54}

\lecturefigure{hyung-slide-055.jpg}{Un modelo de lenguaje general debe integrar conocimiento entre idiomas, lo que debilita la motivación de los parámetros independientes}{Hyung Won Chung official deck, p.~55}

Al leer la segunda diapositiva conviene notar que la tarea ya cambió. El preentrenamiento moderno no solo aprende traducción inglés-alemán, sino que debe formar un conocimiento compartido del mundo a partir de texto multilingüe; un mismo hecho puede expresarse en distintos idiomas. Si la entrada y el objetivo quedan aislados de forma permanente, al modelo puede costarle más reutilizar las representaciones. Por otro lado, la especialización multilingüe todavía puede tener valor; por eso esta diapositiva plantea un scale-era counterargument, no que los parámetros independientes sean absolutamente inútiles.

\teachervoice{La comparación histórica de 00:57:51--00:59:09 es muy concreta: cuando el objetivo es solo traducir, separar los parámetros de los distintos idiomas es «muy natural»; cuando el objetivo pasa a ser el aprendizaje de conocimiento general, el ponente prefiere fusionar el conocimiento entre idiomas. La evaluación de un inductive bias cambia junto con el objective.}

\subsection{FLAN: entradas largas y objetivos cortos dan una ventaja inesperada a encoder-decoder}

El instruction tuning ofrece un experimento natural más fino. El mismo equipo hizo fine-tuning sobre PaLM (decoder-only) y T5 (encoder-decoder); aunque dedicó la mayor parte del tiempo a optimizar PaLM, la ganancia de T5 fue mayor. El ponente no atribuye el resultado a una superioridad general de la arquitectura, sino que vuelve a la distribución de longitudes de los datos para buscar la correspondencia.

\lecturefigure{hyung-slide-056.jpg}{El instruction finetuning unifica una gran cantidad de tareas académicas como instrucciones en lenguaje natural}{Hyung Won Chung official deck, p.~56}

\lecturefigure{hyung-slide-057.jpg}{En el experimento FLAN, la ganancia de fine-tuning de encoder-decoder es mayor}{Hyung Won Chung official deck, p.~57}

\begin{knowledgebox}{Lectura de la figura: se comparan ganancias, no qué modelo gana en términos absolutos}
La tabla destaca el improvement antes y después del instruction tuning. No puede concluirse, mirando solo el recuadro rojo, que T5 supera a PaLM en todas las tareas y a todas las escalas; la conclusión de la clase es que la estructura de encoder-decoder coincide especialmente bien con ese conjunto de datos académicos, por lo que aprovecha mejor la señal del fine-tuning.
\end{knowledgebox}

\lecturefigure{hyung-slide-058.jpg}{Los datos académicos suelen tener una distribución de longitudes con entradas largas y objetivos cortos}{Hyung Won Chung official deck, p.~58}

Si la variable aleatoria de la longitud de la entrada es $L_x$ y la longitud del objetivo es $L_y$, este conjunto de datos cumple aproximadamente
\[
\mathbb{E}[L_x]\gg \mathbb{E}[L_y].
\]
Aquí $\mathbb{E}[L_x]$ y $\mathbb{E}[L_y]$ son, respectivamente, la longitud promedio de la entrada y del objetivo. Los enunciados largos y las etiquetas cortas hacen que las distribuciones de ambos lados sean claramente distintas, de modo que los parámetros independientes del encoder/decoder pueden especializarse; sin embargo, la generación de textos largos y el chat de varios turnos reducen esta asimetría o incluso la invierten.

\teachervoice{En 00:59:46--01:01:10, Hyung Won dice que el equipo dedicó el 99\% del tiempo a optimizar PaLM y, sin embargo, T5 obtuvo una ganancia mayor en pocos días; su explicación es que el sesgo de longitud de los datos y la arquitectura «coincidieron por casualidad». Este teacher voice convierte el resultado de una comparación de marcas en un análisis de correspondencia de distribuciones.}

\begin{warningbox}{Que un academic benchmark sea fácil de calificar moldea la distribución de entrenamiento}
Los objetivos largos son difíciles de anotar manualmente y de calificar de forma automática, por lo que los datos académicos se inclinan hacia entradas largas y respuestas cortas. Que algo sea fácil de calificar no equivale a que tenga valor para el usuario: la escritura creativa, el código, el análisis y el diálogo suelen requerir salidas largas. Si la arquitectura del modelo se optimiza solo para la distribución del benchmark, la ventaja estructural original puede desaparecer cuando cambien las tareas de despliegue.
\end{warningbox}

\subsection{Resumen de la sección}

La traducción automática y FLAN muestran por qué el sesgo inductivo puede producir beneficios reales en una época determinada: dirige la capacidad limitada hacia donde más lo necesita la distribución de los datos. Pero el preentrenamiento general, la generación larga y el diálogo de varios turnos cambiaron la relación entre entrada y objetivo, lo que obligó a los investigadores a volver a probar la separación de parámetros. La historia no sirve para burlarse de los diseños antiguos, sino para identificar las condiciones en las que esos diseños son válidos.
\section{Granularidad de la representación, Bidirectionality y caché en conversaciones de varios turnos}

La sección anterior analizó la separación de parámetros; esta sección aborda las dos estructuras restantes: que el decoder solo lee la última capa del encoder y que el encoder usa bidirectional attention. Ambas mejoraron en su momento la calidad en las tareas, pero a medida que las redes se vuelven más profundas y las conversaciones más largas, los problemas de granularidad de la información y de cálculo repetido se vuelven notorios. Aquí, por primera vez, se colocan la architecture choice y el serving cost en una misma figura.

\subsection{¿Leer solo la última encoder layer crea un cuello de botella de información?}

Esta subsección se centra en el punto de lectura de la cross-attention. En las redes profundas, distintas capas suelen representar información de distinta granularidad; si incluso la primera capa del decoder solo puede acceder a la última capa del encoder, la información formal de las capas bajas podría quedar excesivamente comprimida. En los modelos actuales de algunas decenas de capas el efecto quizá sea pequeño, pero en redes extremadamente profundas este desajuste podría amplificarse.

\lecturefigure{hyung-slide-059.jpg}{Segunda suposición: el objetivo solo necesita leer la entrada ya completamente codificada}{Hyung Won Chung official deck, p.~59}

\lecturefigure{hyung-slide-060.jpg}{La cross-attention estándar solo lee las activaciones de la última capa del encoder}{Hyung Won Chung official deck, p.~60}

\lecturefigure{hyung-slide-061.jpg}{En las redes profundas, las capas inferiores y superiores suelen codificar información de distinta granularidad}{Hyung Won Chung official deck, p.~61}

\begin{knowledgebox}{Lectura de la figura: la representation hierarchy es una heurística, no una tabla semántica fija}
En las redes de visión se suele usar «bordes en las capas bajas, objetos en las capas altas» para explicar la representación jerárquica; en los modelos de lenguaje también puede aparecer una tendencia de la forma local a la semántica abstracta. Sin embargo, el significado de cada capa cambia con el objetivo de entrenamiento, las conexiones residuales y la profundidad. Lo que la figura realmente respalda es que distintas capas pueden aportar información complementaria; por lo tanto, exponer solo la última capa es una decisión estructural que debe ponerse a prueba.
\end{knowledgebox}

\teachervoice{01:02:43--01:03:25, Hyung Won dice que en los T5 de unas 24 capas que él ha usado, que la layer 1 lea la layer 24 parece no causar problemas; pero si en el futuro hubiera diez o mil veces más profundidad, él «no se sentiría del todo cómodo». Es una pregunta típica de scale extrapolation, no una falla ya observada.}

\subsection{Qué contrafactual es la Layer-aligned connection}

La subsección anterior señaló el problema de granularidad; esta construye la alternativa mínima: que la capa $\ell$ del decoder lea la capa $\ell$ del encoder. Así, cada etapa recibe una representación de entrada de profundidad similar, lo que además se acerca más a la token interaction de la misma capa en un decoder-only. El valor del contrafactual está en que permite hacer experimentos controlados, aunque al final no se adopte.

\lecturefigure{hyung-slide-062.jpg}{Tercera suposición: dentro de la entrada se necesita una bidirectional interaction all-to-all}{Hyung Won Chung official deck, p.~62}

\lecturefigure{hyung-slide-063.jpg}{Arquitectura alternativa con layer-aligned connection y causal input}{Hyung Won Chung official deck, p.~63}

\begin{importantbox}{El producto correcto del análisis estructural es un ablation axis}
De H059--H063 se obtienen dos ejes experimentales independientes: final-layer versus layer-aligned cross-attention, y bidirectional versus causal input mask. Si se cambian ambos a la vez, la diferencia de desempeño no puede atribuirse a ninguno; la derivación en cuatro pasos del curso sirve precisamente para generar ablations que puedan escalarse por separado.
\end{importantbox}

\subsection{Beneficio en calidad y costo de sistema de la Bidirectionality}

El bidirectional encoder fue muy importante en la época de BERT para tareas de comprensión difíciles, porque cada token puede aprovechar a la vez el contexto izquierdo y el derecho. Sin embargo, en la instruction tuning a gran escala, la anecdotal experience del ponente es que la diferencia entre bidireccional y unidireccional es pequeña; al mismo tiempo, las conversaciones de varios turnos exhibieron un costo de sistema evidente: un mensaje nuevo cambia el «futuro» de la entrada anterior, y la representación bidireccional tiene que volver a codificarse.

\lecturefigure{hyung-slide-064.jpg}{A escala suficiente, el beneficio de la bidirectionality puede reducirse, pero aumenta el costo en las conversaciones de varios turnos}{Hyung Won Chung official deck, p.~64}

\teachervoice{01:03:35--01:04:23, Hyung Won marca explícitamente la conclusión como highly anecdotal: en Flan-2, la diferencia entre fine-tuning bidireccional y unidireccional es pequeña. Las notas conservan este nivel de evidencia; no debe reescribirse como «se ha demostrado que la bidirectional attention es inútil».}

\subsection{Por qué la causal attention puede reutilizar la KV cache}

Esta subsección convierte la figura en una cuenta de cómputo. En un causal model, la representación de los tokens anteriores no depende de los tokens nuevos futuros, por lo que sus key/value calculados en el turno anterior pueden guardarse en caché; al añadir un mensaje solo se calcula la parte nueva. En un bidirectional encoder, los tokens anteriores pueden leer los tokens nuevos, de modo que después de añadir algo las representaciones anteriores dejan de ser válidas y, por lo general, hay que volver a codificar toda la entrada.

\lecturefigure{hyung-slide-066.jpg}{Comparación entre la recodificación bidirectional y la cache reuse unidirectional en conversaciones de varios turnos}{Hyung Won Chung official deck, p.~66}

Sea $n$ la longitud del contexto existente y $m$ la longitud añadida. Ignorando constantes y el MLP, el costo de rehacer la full attention es aproximadamente
\[
O\bigl((n+m)^2\bigr),
\]
mientras que la causal KV cache (key-value cache, que guarda los tensores key/value de cada capa para los tokens anteriores) solo necesita calcular, para las nuevas query, la attention sobre la caché anterior y sobre los tokens añadidos, con un costo incremental aproximado de
\[
O(nm+m^2).
\]
Aquí $n$ es la longitud del contexto anterior y $m$ es el número de tokens añadidos. La caché reduce el cálculo repetido, pero aumenta el uso de memoria de la GPU; el servicio con contextos largos todavía requiere gestionar la cache placement, el batching y la eviction.

\begin{knowledgebox}{Lectura de la figura: la región azul reutilizable corresponde a la invariancia causal}
En el caso unidirectional, los tokens anteriores no pueden ver el futuro, por lo que la llegada de un mensaje nuevo no cambia los key/value anteriores; en el caso bidirectional, la representación de ``How are you?'' cambia porque después aparece ``Bad''. La ventaja de eficiencia del sistema proviene de la dirección de las dependencias en el grafo de cómputo, no del nombre decoder-only en sí.
\end{knowledgebox}

\teachervoice{La explicación en clase de 01:04:31--01:05:36 conecta la calidad del modelo con el serving: una estructura que en 2018 mejoró los benchmarks puede generar costos de recodificación en la forma de producto de varios turnos. La architecture trajectory no la determina solo la accuracy, sino también el modo de interacción.}

\subsection{Conclusión y Q\&A: el siguiente cuello de botella estructural podría ser el objective}

Al terminar el análisis de arquitectura, el ponente lleva el método de vuelta a la investigación en general. Las estructuras históricas no son reliquias irrelevantes; hacer este análisis repetidamente entrena al investigador para identificar las suposiciones implícitas del problema actual y para preguntarse si pueden sustituirse por un método más general acompañado de escala. La sesión de Q\&A plantea además que el cuello de botella más fuerte en la actualidad quizá no esté en el Transformer block, sino en el objetivo de aprendizaje.

\lecturefigure{hyung-slide-067.jpg}{Conclusión: identificar la fuerza dominante y analizar la estructura adicional desde una perspectiva de scaling}{Hyung Won Chung official deck, p.~67}

\teachervoice{El cierre de 01:06:06--01:06:42 pide a los estudiantes volver a sus propios problemas y preguntarse «qué suposiciones hay que revisar, por qué, y si pueden cambiarse por un mecanismo más general y hacer scale up». Esto es más importante que recordar que ganó el decoder-only, porque el mismo método también sirve para MoE, multimodality, memory y learning objective.}

El entrenamiento por máxima verosimilitud (maximum-likelihood estimation, MLE) toma el objetivo observado como la continuation correcta en las muestras de entrenamiento:
\[
\theta^{\star}=\arg\max_{\theta}\sum_{(x,y)\in\mathcal{D}}\log p_{\theta}(y\mid x).
\]
Aquí $\mathcal{D}$ es el conjunto de datos de pares entrada-objetivo, $x$ es la condición, $y$ es el objetivo observado y $\theta$ son los parámetros del modelo. En traducción o clasificación, una sola respuesta de referencia suele proporcionar una supervisión eficaz; en tareas abiertas como «escribe un poema», tratar un único $y$ como la única señal positiva impone una suposición estructural muy fuerte.

\begin{knowledgebox}{Qué papel desempeña aquí el RLHF}
El RLHF (reinforcement learning from human feedback) primero entrena un reward model a partir de preferencias humanas y luego usa esa recompensa para optimizar la política. Hyung Won lo considera un ejemplo de «aprender la objective function»: en lugar de solo copiar un único objetivo, se hace que el modelo aprenda la calidad relativa de varias respuestas. Al mismo tiempo, enfatiza que el RLHF tampoco es lo bastante scalable, así que esto solo demuestra una dirección; no es la solución definitiva.
\end{knowledgebox}

\teachervoice{01:12:39--01:14:34, Hyung Won dice que no cree que la architecture sea el principal bottleneck actual; le inquieta más la señal de enseñanza «solo este target es correcto» que el MLE impone a las tareas abiertas. Esta sesión de Q\&A amplía toda la clase: el inductive bias que hay que retirar también puede estar oculto en el objective.}

\begin{warningbox}{La Compute trajectory no es una promesa incondicional}
A partir de 01:15:17, el ponente califica la Moore's law de red herring y enfatiza que la variable real es la compute availability; enumera las GPU, la baja precisión y la especialización del hardware, y también reconoce que la energía y la física pueden convertirse en cuellos de botella. Un plan de investigación debe hacer análisis con varios escenarios, en lugar de extrapolar mecánicamente la tendencia exponencial pasada como una garantía para el futuro.
\end{warningbox}

\subsection{Resumen de la sección}

La jerarquía de representaciones nos recuerda revisar el final-layer bottleneck; las conversaciones de varios turnos nos recuerdan considerar la attention direction junto con el costo de la caché; la sesión de Q\&A, a su vez, extiende el análisis estructural al MLE y al RLHF. El método común sigue siendo: identificar las suposiciones que el sistema actual acepta por defecto, explicar por qué fueron útiles en el regime anterior y luego usar los cambios de escala, de datos y de interacción para comprobar si todavía vale la pena conservarlas.
\section{Síntesis y ampliación}

Las dos conferencias convergen finalmente en una misma disciplina de investigación: no mistificar las capacidades, las arquitecturas ni las tendencias; primero leer los datos y las gráficas, descomponer las tareas y los supuestos ocultos, y después buscar los puntos de cruce a lo largo del eje de escala. Jason ofrece una descomposición microscópica del comportamiento del modelo y Hyung Won, una descomposición macroscópica de la evolución de las arquitecturas; ambos exigen al investigador distinguir entre observación, explicación y extrapolación.

\subsection{Doce conclusiones centrales}

Esta sección condensa las extensas notas en una lista de verificación reutilizable. Cada conclusión conserva sus condiciones de aplicación, para evitar convertir un juicio de clase de 2024 en un hecho de 2026 sin límites.

\begin{enumerate}
  \item Revisar los datos a mano entrena la representación que el propio investigador tiene de la tarea y no puede sustituirse por completo con un benchmark dashboard.
  \item La next-token prediction reúne, mediante continuations reales, una enorme cantidad de latent-task supervision.
  \item La cross-entropy global es una mezcla ponderada de distintas distribuciones de tokens y dificultades de tarea.
  \item Que la loss global sea suave no exige que cada benchmark metric lo sea.
  \item Una emergent-looking curve puede contener a la vez un cambio real de capacidad, una puntuación por umbral y un muestreo disperso.
  \item El inverse/U-shaped scaling suele indicar estrategias en competencia; conviene descomponerlo en subtareas y medirlas.
  \item Una mejora en un solo punto no muestra si el método se satura, si desplaza la curva o si cambia su pendiente.
  \item El cómputo efectivo cada vez más barato es una tendencia histórica sólida, pero no una garantía ilimitada.
  \item El inductive bias mejora la eficiencia en el regime actual, pero puede convertirse en un cuello de botella a mayor escala.
  \item La diferencia entre encoder-decoder y decoder-only puede descomponerse en cuatro hipótesis locales que admiten ablation.
  \item La translation y las tareas de entrada larga con objetivo corto explican por qué ciertas estructuras fueron eficaces en su momento.
  \item La próxima estructura que habrá que reexaminar puede estar en el objective, en la interfaz de datos o en la interacción del sistema, y no solo en el architecture block.
\end{enumerate}

\subsection{Una tabla que conecta a los dos ponentes}

Para situar el comportamiento microscópico y la arquitectura macroscópica en un mismo marco, la siguiente tabla alinea las dos conferencias según «objeto de observación, unidad de descomposición, eje de escala, riesgo y acción». La tabla no presenta conclusiones nuevas; ayuda al lector a trasladar el método.

\begin{table}[H]
\centering
\begin{tabularx}{\textwidth}{L{0.18\textwidth} X X}
\toprule
Dimensión & Jason Wei & Hyung Won Chung \\
\midrule
Objeto de observación & token, prompt, task metric, loss curve & arquitectura, tareas históricas, compute trend, system cost \\
Unidad de descomposición & latent task, subskill, metric threshold & parameter sharing, layer connection, attention mask, objective \\
Eje de escala & model/data/compute budget & compute availability, network depth, interaction length \\
Riesgo principal & Tomar la puntuación agregada por un mecanismo y un salto abrupto por una capacidad misteriosa & Volver permanentes los atajos históricos y tomar los benchmarks antiguos por las tareas futuras \\
Acción recomendada & Mirar los datos, descomponer las tareas, trazar scaling curves con varios puntos & Identificar la fuerza dominante, registrar los sesgos, hacer ablations por sustracción \\
\bottomrule
\end{tabularx}
\caption{Correspondencia entre las dos conferencias dentro de un mismo marco guiado por la evidencia.}
\end{table}

\subsection{Plantilla de investigación para la práctica}

El lector puede aplicar el método del curso directamente a una idea nueva. Primero hay que establecer qué observación es un hecho y después formular al menos dos explicaciones en competencia; a continuación, se eligen los puntos de presupuesto, las subtareas o las ablations de arquitectura que permitan distinguir entre esas explicaciones. Si el resultado solo se sostiene en un punto, no debe afirmarse que cambia la scaling trajectory.

\begin{importantbox}{Plantilla de siete pasos: del aula al experimento}
\begin{enumerate}
  \item Muestrear y revisar a mano datos reales de entrenamiento o de evaluación;
  \item Definir una metric continua y una tarea discreta del usuario, para no quedarse con una sola puntuación;
  \item Enumerar las tareas latentes, las estrategias y las hipótesis estructurales;
  \item Ejecutar la baseline y la intervention en varios puntos de presupuesto;
  \item Hacer ablations de una sola variable sobre las estructuras clave;
  \item Reportar los intervalos de ajuste, los errores, los casos de falla y la distribución de los datos;
  \item Indicar con claridad si la conclusión es una observación, una hipótesis de mecanismo, un juicio personal o una extrapolación.
\end{enumerate}
\end{importantbox}

\subsection{Preguntas de autoevaluación}

Estas preguntas sirven para comprobar si se domina de verdad el razonamiento de las conferencias y no solo se recuerdan las conclusiones. Al responder, conviene citar fórmulas, variables de las figuras o experimentos contrafácticos.

\begin{enumerate}
  \item ¿Por qué el next-token objective puede verse como massive multi-task learning? ¿Y por qué no garantiza que todas las tareas se aprendan de forma equilibrada?
  \item Cuando la loss global desciende de manera suave, ¿por qué la exact-match accuracy puede subir de golpe?
  \item ¿Cómo se explica el U-shaped scaling mediante repeat, quote repair e instruction following?
  \item ¿Por qué no basta con que un método de investigación lleve la delantera en un solo punto de presupuesto para afirmar que es más scalable?
  \item ¿En qué condiciones el inductive bias ayuda y en cuáles se convierte en un cuello de botella?
  \item ¿Qué hipótesis elimina cada uno de los cuatro pasos de la transformación de encoder-decoder a decoder-only?
  \item ¿Por qué los datos long-input/short-target podrían ser especialmente adecuados para encoder-decoder?
  \item ¿Por qué la bidirectional attention aumenta el cómputo repetido en los diálogos de varios turnos?
  \item ¿Qué supuesto estructural impone la MLE a la generación abierta? ¿Y qué problema sigue sin resolver el RLHF?
  \item Si en el futuro la compute trend cambia por restricciones energéticas, ¿cómo debería actualizarse el método de análisis de Hyung Won?
\end{enumerate}

\subsection{Lecturas adicionales}

Esta sección enumera solo materiales originales directamente relacionados con la argumentación de la clase. Al leerlos, se recomienda reproducir los ejes y los supuestos de las figuras en lugar de limitarse a extraer las conclusiones.

\begin{itemize}
  \item Kaplan et al., \emph{Scaling Laws for Neural Language Models}: relación empírica de ley de potencias de la loss global.
  \item Wei et al., \emph{Emergent Abilities of Large Language Models}: curvas a nivel de tarea y emergent behavior.
  \item Schaeffer et al., \emph{Are Emergent Abilities of Large Language Models a Mirage?}: efecto de la metric choice en la forma de las curvas.
  \item Sutton, \emph{The Bitter Lesson}: relación a largo plazo entre los métodos generales, el cómputo y la estructura diseñada a mano.
  \item Vaswani et al., \emph{Attention Is All You Need}: el Transformer encoder-decoder original.
  \item Chung et al., \emph{Scaling Instruction-Finetuned Language Models}: evidencia de instruction tuning de la serie FLAN.
\end{itemize}

\end{document}
